\chapter{Competitor Analysis}

\section{The Competitive Field}

The AMR field divides into a few groups. In warehousing and manufacturing the
established players are MiR (now part of Teradyne), Geek+, Locus Robotics, Hai
Robotics, Exotec and Fetch/Zebra, with Robotnik serving the research and
mobile-manipulation niche \cite{standardbots_amr,fdatabot_amr2026,robotnik_kairos}.
Industrial inspection is led by Boston Dynamics (Spot) alongside rugged outdoor
specialists such as Capra and DEEP Robotics \cite{standardbots_amr,generationrobots_amr}.
Security is served by Knightscope, Cobalt Robotics and SMP Robotics, plus the
Tunisian incumbent Enova Robotics \cite{mordor_security,coherent_security}.
Healthcare and hospitality are dominated by purpose-built devices (Aethon TUG,
Savioke), and retail by scanning robots such as Bossa Nova
\cite{mir_healthcare,nordmodules_retail}. A structured comparison is given in
Table~\ref{tab:competitors}.

\section{Cross-Competitor Patterns}

Four patterns stand out across the field.

\textbf{1. Modularity already exists at the top of the market.} The
``one base, many missions'' idea is not new -- MiR's base already accepts shelf
carriers, hooks, pallet lifts and a robotic arm. The differentiator therefore
cannot be modularity in the abstract; it must be price, local support, and ease
of deployment on \emph{unmodified} facilities. The genuine white space is
narrower than it first appears: low-cost modularity for SMEs and legacy
buildings.

\textbf{2. Shelf access splits into two very different things, and the harder
one is crowded.} Vertical reach is dominated by well-capitalised specialists
(Exotec, Hai Robotics, Geek+, Brightpick), most of which require dedicated,
engineered racking rather than retrofitting a customer's existing shelves. The
retrofit-onto-existing-shelves idea is genuinely differentiated but is also the
riskiest engineering claim in the platform and competes against firms with
nine-figure funding -- a reason to defer it, not to lead with it.

\textbf{3. Unit economics are unforgiving, and capex-light models win.} Even
category leaders have struggled financially (Locus Robotics went through
Chapter~11 in 2023), and survivors increasingly sell Robots-as-a-Service
\cite{builtin_raas}. For a Tunisian entrant with limited capital and
conservative buyers, a subscription/lease model is effectively the only way to
clear the procurement barrier.

\textbf{4. Every competitor is calibrated for high-purchasing-power markets.}
This is the decisive pattern for Tunisia: all of the players above price and
package for markets where the robot-cost-to-labor-cost ratio favours
automation. None has a business model adapted to a market like Tunisia, and the
whole competitive set is thin on the ground in North Africa -- global players
reach the GCC through integrators but have essentially no local manufacturing,
support, or French/Arabic presence locally. That gap is precisely where a local
team can win.

\vspace*{0.3cm}
    \begin{center}
    \renewcommand{\arraystretch}{1.6}
    
    \makebox[\textwidth][c]{%
    \begin{tabular}{|>{\centering\arraybackslash}m{2.5cm}|
                    >{\centering\arraybackslash}m{2cm}|
                    >{\centering\arraybackslash}m{2.5cm}|
                    >{\centering\arraybackslash}m{3cm}|
                    >{\centering\arraybackslash}m{1.5cm}|
                    >{\centering\arraybackslash}m{1.5cm}|
                    >{\centering\arraybackslash}m{2cm}|
                    >{\centering\arraybackslash}m{0.5cm}|}
    
    \hline
    \rowcolor{gray!15}
     \textbf{Company} &\textbf{Product }&\textbf{Sector  }  & \textbf{Key features} & \textbf{Pricing } & \textbf{Strengths } & \textbf{Weaknesses}& \textbf{Source} \\
    \hline
    
    \textbf{(Teradyne)}& MiR250 / 600 / 1350&  Warehousing, Mfg.&Interchangeable top modules; arm option (MC600)  & USD 30--45k& Mature, large install base &Proprietary software lock-in & \cite{mir250} \\
    \hline
    \textbf{Robotnik}& RB-KAIROS+& Mfg., R\&D& ROS2 base, plug-and-play arm &From EUR 50k &Open, modular & High entry price  & \cite{robotnik_kairos}\\
    \hline
    \textbf{Geek+}& P-series, RoboShuttle&Warehousing&Goods-to-person, sortation, shuttle &  Project-based  & Huge fleet scale & System-level, not one robot & \cite{fdatabot_amr2026}\\
    \hline
    \textbf{Locus Robotics}& LocusBots & Warehousing, 3PL & Collaborative picking, RaaS & RaaS subscription & Capex-light model & Single workflow; financial stress & \cite{standardbots_amr,builtin_raas} \\
    \hline
    \textbf{Hai Robotics}& HaiPick / Climb & Warehousing (totes) & Tote/case ACR, vertical reach & Project-based & Vertical density & Needs engineered racking & \cite{standardbots_amr} \\
    \hline
       \textbf{Exotec}&Skypod & Warehousing (high-density) & Shelf-climbing goods-to-person & Project, multi-million & High throughput & System-level only, costly & \cite{fdatabot_amr2026} \\
    \hline       \textbf{Boston Dynamics}&Spot & Inspection & Legged; camera/sensor mounts & USD 75k+ & Best-in-class mobility & Very high cost and training & \cite{standardbots_amr} \\
    \hline       \textbf{Capra Robotics}&Capra Hircus & Outdoor inspection/patrol & Fully modular, rugged & Quote-based & Outdoor-rugged & Narrow niche & \cite{generationrobots_amr} \\
    \hline

      \textbf{Fetch / Zebra}& Freight AMRs & Warehouse transport & Transport + data & Project / fleet & Zebra channel reach & Transport-focused only & \cite{standardbots_amr} \\
    \hline
      \textbf{Aethon}& TUG & Healthcare, hospitality & Secure cart delivery & Quote-based subscription & Proven in hospitals & Single-purpose, not reconfigurable & \cite{researchgate_hospital} \\
    \hline
      \textbf{Cobalt Robotics}&Indoor patrol & Security & Patrol + teleoperation & USD 3.5--6.5k/mo & Indoor security & Security-only & \cite{coherent_security} \\
    \hline
   
  \end{tabular}%
    }
    \end{center}
       \begin{center}
    \renewcommand{\arraystretch}{1.6}
    
    \makebox[\textwidth][c]{%
    \begin{tabular}{|>{\centering\arraybackslash}m{2.5cm}|
                    >{\centering\arraybackslash}m{2cm}|
                    >{\centering\arraybackslash}m{2cm}|
                    >{\centering\arraybackslash}m{3cm}|
                    >{\centering\arraybackslash}m{1.5cm}|
                    >{\centering\arraybackslash}m{1.5cm}|
                    >{\centering\arraybackslash}m{2cm}|
                    >{\centering\arraybackslash}m{0.5cm}|}
    
    \hline
   \textbf{SMP Robotics}& Outdoor patrol & Security & Additional sensors & USD 120--320k & Large-area patrol & Extremely high cost & \cite{mordor_security} \\
    \hline

      \textbf{Bossa Nova}& Shelf scanner & Retail & Scanner only & Undisclosed & Shelf analytics & Scanning only; Walmart ended use & \cite{nordmodules_retail,dudarenok_retail_china} \\
    \hline
     \textbf{Enova Robotics (TN)}&PGuard & Security, inspection & Custom design + AI; locally built & Undisclosed & Local capability; exported & UGV/security-focused & TBD (local) \\
    \hline



    
    \end{tabular}%
    }
    \end{center}







